
Let $A\ingroundset\integers$ with $\average(A)=\mu_A\in\integers$ and 
$\nofdroplets{A}=n_A$ be an arbitrary configuration (that is, a multiset of integers).
We define $\min(A)$ and $\max(A)$ as the lowest and highest concentrations in $A$, respectively.
We also define the diameter of $A$ as $\diameter(A) = \max(A) - \min(A)$.
Recall that the size of $A$ is defined as $s(A)=\sum_{a\in A} \log(|a|+2)$. 
We denote by $A_{even}$ and $A_{odd}$ two disjoint multisets containing
all even and odd concentrations in $A$, respectively.
We use notation $\pi\in\braced{even,odd}$ for a parity value, with
$\barpi\in\braced{even,odd}$ standing for the opposite parity (that is $\barpi\neq\pi$),
and we use these for subscripts, as in  $A_\pi$ and $A_{\barpi}$.   
In the contexts when $A$ is subjected to some mixing sequence,
we will use notation $\potential_{A,0}$ as the initial value of $\potential(A)$,
before any mixing has been performed on $A$.

The following observations show that repeatedly mixing two droplets with concentrations 
that are sufficiently far apart eventually decreases $\potential(A)$ at least by a factor of $\half$.

%%%%%%%%%%%%%

\begin{observation}\label{obs: far-appart mix}
Let $x,y\in A$ be two different concentrations with the same parity, and let $A'$ be
obtained from $A$ by mixing $x$ and $y$. Then:
%
\begin{description}
	
	\item{(a)} $\potential(A') < \potential(A)$.

	\item{(b)} If $\gamma\in\posintegers$ is a constant and
	$|x-y|\geq\diameter(A)/\gamma$, then
	$\potential(A') \;\leq\; \big(1 - \frac{1}{2\gamma^2n_A} \big)\potential(A)$.
	
\end{description}

\end{observation}

\begin{proof}
Without loss of generality we can assume that $\mu_A= 0$, so that $\potential(A)=\sum_{a\in A}a^2$.
Note that we then also have $\average(A')=\mu_{A'} = 0$.
Part~(a) now follows by simple calculation:
%
	\begin{equation*}
			\potential(A) - \potential(A')
			\;=\; x^2+y^2-2\Big(\frac{x+y}{2}\Big)^2 
			\;=\; \frac{(x-y)^2}{2} > 0.
	\end{equation*}

\noindent
To prove~(b), we use the above calculation and inequality $\potential(A)\le n_A \diameter(A)^2$
(that follows directly from the definition of $\potential(A)$):
%
	\begin{equation*}
			\potential(A) - \potential(A')
			\;=\; \frac{(x-y)^2}{2} 
			\;\geq\; \frac{\diameter(A)^2}{2\gamma^2} 
			\;\geq\; \frac{\potential(A)}{2\gamma^2n_A},
	\end{equation*}

\noindent
completing the proof.
\end{proof}


%%%%%%%%%%%%%

\begin{observation}\label{obs: far-appart constant mix}
	Let $\gamma\in\posintegers$ be a constant.   
	Let $A'$ be obtained from $A$ by a sequence of $2\gamma^2n_A$ mixing operations,
	each involving droplets $x,y\in A$ that satisfy $|x-y|\geq\diameter(A)/\gamma$.  
	Then $\potential(A') \le \potential(A)/2$.
\end{observation}

\begin{proof}
	Applying Observation~\ref{obs: far-appart mix}(b) repeatedly, for each of the
	$2\gamma^2n_A$ mixing steps, we obtain
	%
	\begin{equation*}
		\potential(A')
		\leq \potential(A)\Big(1 - \frac{1}{2\gamma^2n_A}\Big)^{2\gamma^2n_A}
		\leq \frac{\potential(A)}{2},
	\end{equation*}
	%
	where the last inequality holds because $(1-1/k)^k \leq 1/2$ for $k\geq 2$.
\end{proof}


%%%%%%%%%%%%%

We state two more observations that will be used later in the proof, typically
without an explicit reference.

\begin{observation}\label{obs: logPotential <= 2s(A)}
	$\log{\potential(A)}\leq 2s(A)$.
\end{observation}

\begin{proof}
	Let $\psi_A(x) =  \sum_{a\in A}(a-x)^2$. By calculus, we have that
	$\psi_A(x)$ is minimized for $x = \mu_A$, and thus
	%
	\begin{equation*}
		 \potential(A) \;=\; \psi_A(\mu_A)
						\;\le\; \psi_A(0)
						\; = \;  \sum_{a\in A}a^2
						\;\le    \Big(\,\sum_{a\in A} |a| \,\Big)^2.
	\end{equation*}   
	%
	Therefore 
	%
	\begin{align*}
		\log\potential(A) &\leq\; \log\Big(\,\sum_{a\in A} |a| \,\Big)^2
		\\
		&\leq\; 2\log\prod_{a\in A}\big(|a|+2\big)
		\;=\; 2\sum_{a\in A}\log(|a|+2)
		\;=\; 2s(A),
	\end{align*}	
	%
completing the proof.
\end{proof}

\begin{observation}\label{obs: psi(A') <= psi(A) for A' subseteq A}
	Let $A'\subseteq A$. Then $\potential(A')\leq\potential(A)$.
\end{observation}      

\begin{proof}  
Using the properties of the 
function $\psi_A(x)$ from the proof of Observation~\ref{obs: logPotential <= 2s(A)},
we have
%
\begin{equation*}	 
	 \potential(A') \;=\; \psi_{A'}(\mu_{A'})
	 				\;\leq\; \psi_{A'}(\mu_{A})   
	 				\;\leq\; \psi_{A}(\mu_{A})  
	   				\;=\; \potential(A),
\end{equation*}	
%
where the second inequality follows from $A'\subseteq A$.  	
\end{proof}





